\par\Needspace{6\baselineskip}\medskip
\begin{center}\hypertarget{native-182-article136}{}\pdfbookmark[2]{136.}{article-136}136.\end{center}
First of all, we must join the three heliocentric positions of the Earth on the celestial sphere
$A,A'\;A''$ (Fig. 4) to the three corresponding geocentric positions of the celestial body $B,B',B''$ by
great circles. We must then calculate both the orientation of these great circles with respect
to the ecliptic (assuming that we adopt the ecliptic as the fundamental plane) and the position
of the points $B,B',B''$ on them. Let $\alpha,\alpha',\alpha''$ be the three geocentric longitudes of the celestial body,
$\beta,\beta',\beta''$ its latitudes, and $l,l',l''$ the heliocentric longitudes of the Earth, whose latitudes we set
$=0$ (Articles 117, 72). Let $\gamma,\gamma',\gamma''$ further denote the inclinations to the ecliptic of the great
circles drawn from $A,A',A''$ to $B,B',B''$, respectively. To follow a fixed rule in determining these
inclinations, we shall always measure them with respect to the part of the ecliptic that extends
from the points $A,A',A''$ in the order of the signs of the zodiac. We thus reckon them from $0$
to $360^0$, or equivalently from $0$ to $180^0$ on the northern side and from $0$ to $-180^0$
on the southern side. We denote the arcs $AB,A'B',A''B''$, which we can always take between $0$ and
$180^0$, by $\delta,\delta',\delta''$. To determine $\gamma,\delta$ we thus have the formulas
\[
\begin{aligned}
1)\quad&\operatorname{tang}\gamma=\frac{\operatorname{tang}\beta}{\sin(\alpha-l)}\\[.7ex]
2)\quad&\operatorname{tang}\delta=\frac{\operatorname{tang}(\alpha-l)}{\cos\gamma}\mathpunct{\raisebox{.45ex}{,}}
\end{aligned}
\]
to which, if desired, the following may be added as a check on the calculation
\[
\sin\delta=\frac{\sin\beta}{\sin\gamma}\mathpunct{\raisebox{.45ex}{,}}\qquad
\cos\delta=\cos\beta\cos(\alpha-l).
\]
There are clearly entirely analogous formulas for determining $\gamma',\delta',\gamma'',\delta''$. If we had $\beta=0$ together
with either $\alpha-l=0$ or $=180^0$, i.\,e. if the celestial body were at opposition or conjunction and
at the same time on the ecliptic, $\gamma$ would be indeterminate. We assume that this case occurs
in none of the three observations.

If the equator has been adopted as the fundamental plane in place of the ecliptic, determining
the orientation of the three great circles with respect to the equator requires, besides the
inclinations, the right ascensions of the intersections with the \hypertarget{native-183}{}equator. Besides the
distances of the points $B,B',B''$ from these intersections, we must also calculate the distances of
the points $A,A',A''$ from them. As these calculations depend on the problem treated in Article 110,
we do not pursue the derivation of the formulas here.
